% Generated by scripts/tikz_reviews.py using the compiled book's labels.
\pdfbookmark[0]{Chapter 13}{comparison-chapter-13}
\ReviewFigure{fig-review-ml-pca-variance}
{PCA directions and projected variance}
{figures/ml/variance.png}
{figures/tikz/ml-pca-variance.pdf}
{Principal directions are eigenvectors of the empirical covariance; the chapter writes its eigenvalues as sigma squared.}
{The centered illustrative cloud has exact reflection symmetries in principal coordinates. The ellipse semiaxes are empirical standard deviations with normalization 1/n; variance is their square. A highlighted observation and its perpendicular projection show the coordinate whose squared values are averaged. Unit eigenvectors and standard-deviation lengths remain distinct.}
{ml-pca-variance}
{Complete figure}
{13.2}
\ReviewFigure{fig-review-ml-pca-data-matrix}
{Rows are observations; columns are features}
{figures/ml/pca-maths/pca-maths-1.pdf}
{figures/tikz/ml-pca-data-matrix.pdf}
{The PCA proof centers the data and writes the n-by-p matrix X with observations in its rows.}
{The reconstruction follows the current n-observations, p-features convention. An observation x\_i is a column vector in R\textasciicircum{}p; the corresponding matrix row is its transpose. Centering is explicit, and the covariance is X transpose X divided by n.}
{ml-pca-data-matrix}
{Complete figure}
{13.4}
\ReviewFigure{fig-review-ml-pca-compression-reconstruction}
{Linear compression followed by reconstruction}
{figures/ml/pca-maths/pca-maths-2.pdf}
{figures/tikz/ml-pca-compression-reconstruction.pdf}
{The variational PCA problem compares rank-at-most-k linear reconstructions R S transpose.}
{Both R and S have p rows and k columns. Compression applies S transpose; reconstruction applies R, so the product is a p-by-p map of rank at most k. Orthogonality is not assumed for arbitrary R and S; it is justified later when reducing to orthogonal projections.}
{ml-pca-compression-reconstruction}
{Complete figure}
{13.5}
\ReviewFigure{fig-pca-var-proof}
{The polytope bounding the PCA objective}
{figures/ml/pca-maths/pca-maths-3.pdf}
{figures/tikz/ml-pca-linear-program.pdf}
{The PCA proof bounds the trace using \eqref{eq-variational-pca}, a linear program maximizing a decreasingly weighted sum over 0 <= beta <= 1 and sum beta <= k.}
{The left feasible region is a triangle for p=2, k=1. The right is the unit cube truncated by sum beta <= 2 for p=3, k=2. Its cutting face is shown with nonzero visible area; dashed edges distinguish the hidden cube edges from the removed corner (1,1,1). Red dots mark valid maximizers (1,0) and (1,1,0) for ordered nonnegative eigenvalues, without asserting uniqueness.}
{ml-pca-linear-program}
{Panel 1 of 3}
{13.6}
\ReviewFigure{fig-pca-var-proof}
{Row norms of the rotated basis}
{figures/ml/pca-maths/pca-maths-4.pdf}
{figures/tikz/ml-pca-row-norms.pdf}
{The PCA proof defines B=V transpose S and beta\_i as the squared norm of its ith row.}
{The matrix B has p rows and k orthonormal columns. The vector b\_i is in R\textasciicircum{}k and appears transposed as a row. Its squared row norms sum to k, the squared Frobenius norm of B. These dimensions follow the current proof rather than ambiguous letters in the scan.}
{ml-pca-row-norms}
{Panel 2 of 3}
{13.6}
\ReviewFigure{fig-pca-var-proof}
{Completing the columns to an orthogonal matrix}
{figures/ml/pca-maths/pca-maths-5.pdf}
{figures/tikz/ml-pca-orthogonal-extension.pdf}
{An orthogonal completion proves that every squared row norm beta\_i is at most one.}
{The first k columns are B; the remaining p-k columns form B perpendicular. Every full row of the square orthogonal matrix has norm one. Restricting a row to the first k entries cannot increase its norm. When k=p, the complementary block is empty.}
{ml-pca-orthogonal-extension}
{Panel 3 of 3}
{13.6}
\ReviewFigure{fig-review-ml-voronoi}
{Nearest-centroid assignment and Voronoi cells}
{figures/ml/voronoi.png}
{figures/tikz/ml-voronoi.pdf}
{For fixed centroids, the assignment step of k-means selects the nearest centroid in Euclidean distance.}
{Cell edges are the exact perpendicular bisectors of the displayed centroids. The dashed centroid segment and right-angle marker make that construction explicit. Colored observations lie in their nearest-centroid cells; the common vertex is equidistant from all three centroids. This illustrates assignment, not a claim that these centroids are already the cluster means. Boundary ties require the convention in the text.}
{ml-voronoi}
{Complete figure}
{13.7}
\ReviewFigure{fig-review-ml-joint-model}
{A joint probability model for input and output}
{figures/ml/proba-model.pdf}
{figures/tikz/ml-joint-model.pdf}
{Supervised observations are pairs (x\_i,y\_i) drawn from a joint law pi on the product of input and output spaces.}
{Contours represent an illustrative bivariate Gaussian density, not an asserted fit to the original hand-drawn points. The dashed projections identify the two coordinates of one observation. The general discussion in the chapter is not restricted to Gaussian models.}
{ml-joint-model}
{Complete figure}
{13.10}
\ReviewFigure{fig-bound-regul}
{Conditional expectation as the mean of a slice}
{figures/ml/cond-expect.pdf}
{figures/tikz/ml-conditional-expectation.pdf}
{Under squared loss, the optimal predictor is the conditional expectation of the output given the input.}
{Here X is uniform on [-1.75,1.75] and Y=m(X)+epsilon, with independent centered asymmetric mixture noise: weights 3/4 and 1/4, Gaussian means -0.3 and 0.9, and standard deviation 0.28. Shading encodes joint density; dividing its vertical slice by p\_X(x0) gives the right-hand conditional density. Its mean m(x0) differs from its mode. The theorem does not assume this illustrative noise model.}
{ml-conditional-expectation}
{Complete figure}
{13.11}
\ReviewFigure{fig-review-ml-linear-fit}
{An affine fit and a nonlinear regression function}
{figures/ml/linear-fit.png}
{figures/tikz/ml-linear-fit.pdf}
{Linear regression restricts the predictor to a finite-dimensional model, even when the conditional expectation is nonlinear.}
{The red line is the least-squares affine fit to the actual displayed illustrative points. The dashed curve is the model conditional mean used to construct those points. Residuals are vertical output errors, not orthogonal distances to the line.}
{ml-linear-fit}
{Complete figure}
{13.12}
\ReviewFigure{fig-review-ml-nearest-neighbors}
{Four nearest neighbors of a query point}
{figures/ml/clustering.png}
{figures/tikz/ml-nearest-neighbors.pdf}
{The R-nearest-neighbor rule sorts training points by their Euclidean distance from the query.}
{The arrows identify the four closest observations in increasing distance order. The pale disk is the selected neighborhood; its radius is exactly the distance to the fourth observation on its boundary. Every remaining point is farther away. The panel explains the distance-selection step before voting over class labels.}
{ml-nearest-neighbors}
{Complete figure}
{13.13}
\ReviewFigure{fig-logistic-sharpness}
{Logistic probability and transition width}
{figures/ml/logistic-1d.png}
{figures/tikz/ml-logistic-one-dimension.pdf}
{The probability of class +1 is theta(beta x), with theta(s)=1/(1+exp(-s)).}
{The probability is theta(beta x)=1/(1+exp(-beta x)), with positive slopes 1 and 3. Both curves pass through probability 1/2 at x=0. Colored dots and guides locate the actual 0.1 and 0.9 crossings, and each dimension arrow measures the corresponding width 2 log(9)/beta. The larger slope has one third of the transition width.}
{ml-logistic-one-dimension}
{Panel 1 of 2}
{13.15}
\ReviewFigure{fig-logistic-sharpness}
{The same decision boundary at two parameter norms}
{figures/ml/logistic-2d.png}
{figures/tikz/ml-logistic-two-dimensions.pdf}
{The parameter direction determines a hyperplane normal; its norm controls the width of the logistic probability transition.}
{Both panels use beta=t v with the same unit normal v. The red dashed contours are labeled by their probabilities 0.1 and 0.9; the dimension arrows measure their normal separation, exactly 2 log(9)/t. The boundary at probability 1/2 stays fixed. Parameter scaling controls probability sharpness, not whether a data set is linearly separable.}
{ml-logistic-two-dimensions}
{Panel 2 of 2}
{13.15}
